\section{Discussion}
\label{sec:discussion}

\subsection{Convergent Evidence for LLM Overconfidence}
\label{sec:discussion-convergent}

The most consistent finding across all three studies is that LLMs---particularly
GPT-4---are substantially overconfident in their confidence intervals for quantitative
forecasting. GPT-4's ECE ranges from 18.9\% (Study 3, favorable domain mix) to 43.7\%
(Study 1, equity-only), with a cross-study mean exceeding 35\%. Even Claude, the
best-calibrated model in every study, achieves only 19.7\% ECE in Study 1---far from
the 0\% target. The pattern of overconfidence is robust to domain, horizon, and sample
size, confirming that it is a systematic property of these models rather than an
artifact of any single experimental design.

These ECE values are comparable to the most severe human expert overconfidence
documented in the literature. \citet{BenDavid2013} found that CFOs' 80\% confidence
intervals for annual equity returns achieved only 36--38\% empirical coverage---roughly
equivalent to our Study 1 Claude result (58\% at 80\%). GPT-4's Study 1 performance
(31\% at the 80\% level) is dramatically worse, suggesting that some frontier LLMs may
exhibit human-level or super-human overconfidence in certain forecasting contexts.

\subsection{The Prediction Horizon Effect}
\label{sec:discussion-horizon}

Study 2 reveals a sharp divergence between models as the prediction horizon grows.
Claude and GPT-4 both degrade substantially, consistent with the hypothesis that their
CI widths are anchored to short-horizon volatility estimates and fail to scale
appropriately with forecast uncertainty. This mirrors findings from the human
overconfidence literature, where longer horizons are consistently associated with
greater interval underestimation \citep{Griffin1992}.

Gemini's behavior is more complex. Its near-perfect 1-day calibration (ECE = 1.8\%,
hit@80 = 78.5\%) and recovery to excellent long-horizon calibration (ECE = 2.5\% at
22 days) suggest that its CI widths implicitly scale with realized volatility in a way
that happens to match empirical coverage at the extremes of the horizon range. Whether
this reflects a genuine uncertainty representation or an accidental match between
Gemini's CI scaling function and stock return volatility at these specific horizons
is an open question that warrants longitudinal replication.

\subsection{A Domain-Confidence Inversion}
\label{sec:discussion-domain}

The central finding of Study 3 is a systematic domain-confidence inversion: LLMs are
most overconfident in financial domains (equities, commodities) where actual
predictability is low, and most underconfident in speculative domains (cryptocurrency,
forex) where they appear to apply large epistemic-humility priors. This pattern is
inconsistent with a single domain-agnostic miscalibration mechanism and suggests
models apply domain-level heuristics inherited from forecasting commentary in their
training corpora.

This finding parallels the human expert literature. \citet{Griffin1992} showed that
perceived predictability often exceeds actual predictability in structured,
information-rich domains. LLMs trained on large corpora of equity research and
commodity analysis may exhibit the same pattern---asserting narrow intervals in domains
where training data contains many precise price forecasts, without encoding the base
rate of forecast failure.

The commodity result is particularly striking: $\mu = 0.25$ uniformly across all three
models. This uniformity suggests a shared domain prior in training data rather than a
model-specific artifact. The crypto inversion (Claude $\mu = 14.4$) implies that models
apply a ``crypto is unpredictable'' narrative that dramatically overstates typical
24-hour volatility.

\subsection{Model Differences}
\label{sec:discussion-models}

GPT-4 is consistently the most overconfident model across all three studies. Its Study
1 ECE (43.7\%) is nearly double Claude's (19.7\%) in the same task, and its Study 2
ECE at 22-day horizon (49.9\%) approaches the theoretical maximum for three-level
calibration. Claude is the best-calibrated model in Study 3 (ECE = 5.3\%) and
competitive with Gemini in Studies 1 and 2. Gemini occupies an intermediate position
overall but is distinctive for its horizon robustness in Study 2.

Claude's extreme underconfidence in cryptocurrency ($\mu = 14.4$) represents the
largest positive outlier in the meta-knowledge analysis, indicating that even the
best-calibrated model can exhibit severe directional miscalibration in specific domains.
GPT-4's forex failure (ECE = 44.2\%) is the worst domain-model result in Study 3 and,
as noted, reflects a compound failure of both point-estimate anchoring and CI width
calibration that the MEAD/MAD ratio alone cannot fully capture.

\subsection{Limitations and Future Directions}
\label{sec:discussion-limitations}

Several limitations constrain our conclusions. First, Studies 1 and 3 were each
conducted on a single day, and Study 3's domain-specific results may partly reflect
idiosyncratic market conditions on March 25, 2026. Study 2's backtest design mitigates
this concern for equities but introduces potential look-ahead biases if models' training
data extends to the backtest period. Second, statistical precision is limited in small
Study 3 domains: $n = 10$ for commodities, $n = 8$ for forex. Third, NBA outcomes were
excluded from Study 3 due to API resolution failures. Fourth, models were evaluated
without chain-of-thought or explicit calibration prompting; whether such interventions
produce genuine calibration improvement is an important open question.

Future work should address these limitations by: (i) collecting predictions daily over
extended periods to separate stable calibration properties from single-day volatility;
(ii) extending the domain set to include earnings forecasts, macroeconomic indicators,
and prediction market prices; (iii) evaluating calibration-specific prompting and
post-hoc recalibration (conformal prediction, temperature scaling); and (iv) using
larger $n$ within Study 3 domains to enable robust domain-level inference.


% -----------------------------------------------------------------------
% Conclusion — no separate section file exists; placed here for continuity
% -----------------------------------------------------------------------
\section{Conclusion}
\label{sec:conclusion}

We report three pre-registered studies examining whether frontier LLMs exhibit
systematic overconfidence in quantitative interval estimates, and how miscalibration
varies with prediction horizon and application domain. Across all three studies, all
three models are substantially overconfident, with GPT-4 consistently the worst
(cross-study ECE 19--44\%) and Claude the best. Study 2 reveals that Claude and GPT-4
become increasingly miscalibrated as the prediction horizon grows from 1 to 22 days,
while Gemini maintains near-target coverage at both short and long horizons. Study 3
reveals a domain-confidence inversion: all models overstate certainty in commodity and
equity markets while overstating uncertainty in cryptocurrency and forex, a pattern
consistent with domain-specific miscalibration inherited from training corpora.

These results have direct implications for practice. Users who elicit confidence
intervals from LLMs for financial decision support should not assume that stated
uncertainty reflects calibrated epistemic states. The pattern of miscalibration is most
dangerous in precisely the domains (commodities, equities) and at precisely the horizons
(multi-week) where overconfident forecasting carries the most severe consequences.
Global calibration performance does not transfer across domains or horizons;
domain-stratified and horizon-stratified evaluation is necessary to characterize
model-specific calibration risk.

More broadly, this study demonstrates that the domain-specificity and
horizon-sensitivity of human overconfidence---developed over four decades of experimental
research on human judges---apply with striking regularity to large language models. LLMs
are not immune to the calibration failures of their training data; they may instead
amplify and systematize them at scale.


% -----------------------------------------------------------------------
% Acknowledgements — no separate section file; placed here
% -----------------------------------------------------------------------
\section*{Acknowledgements}

The authors thank the maintainers of the open data APIs used in this study:
yfinance, CoinGecko, Open-Meteo, and ESPN. No external funding was received
for this research. The authors declare no conflicts of interest.
